% 总体架构：横向箭头对应每问的数据、方法和产出。
% 不在此图中表达尚未核实的跨问数值校准关系。
\begingroup
\definecolor{archblue}{HTML}{315A79}
\definecolor{archlight}{HTML}{EDF3F7}
\definecolor{archgray}{HTML}{536572}
\begin{tikzpicture}[
  font=\song\fontsize{10}{13}\selectfont,
  archbox/.style={draw=archblue!55, line width=0.55pt, rounded corners=2pt,
    fill=archlight, align=center, inner sep=7pt, minimum height=1.55cm},
  archdata/.style={archbox, text width=2.55cm},
  archmodel/.style={archbox, text width=5.1cm},
  archoutput/.style={archbox, text width=3.3cm, fill=white},
  archlink/.style={-{Stealth[length=2mm]}, draw=archblue, line width=0.7pt},
  archhead/.style={text=archblue, font=\song\bfseries\fontsize{11}{14}\selectfont}
]
\node[archhead] at (-5.8,0.65) {数据与条件};
\node[archhead] at (-0.8,0.65) {建模任务与方法};
\node[archhead] at (4.8,0.65) {主要产出};

\node[archdata] (a) at (-5.8,-0.75) {附件 A\\质量信号\\配方与验证损失};
\node[archmodel, fill=archblue, text=white] (q1) at (-0.8,-0.75)
  {问题一：数据质量与领域配比\\熵权评价、冲突分析、岭回归};
\node[archoutput] (o1) at (4.8,-0.75)
  {综合质量评分\\配比与损失关系\\推荐配比};

\node[archdata] (b) at (-5.8,-3.0) {附件 B\\规模、训练量\\与质量信息};
\node[archmodel] (q2) at (-0.8,-3.0)
  {问题二：损失与资源的关系\\广义标度律、边际效用与弹性};
\node[archoutput] (o2) at (4.8,-3.0)
  {损失函数与参数\\资源之间的替代关系};

\node[archdata] (c) at (-5.8,-5.25)
  {算力预算\\题给成本设定\\C7 上下文信息};
\node[archmodel] (q3) at (-0.8,-5.25)
  {问题三：算力约束下的配置\\成本建模与约束优化};
\node[archoutput] (o3) at (4.8,-5.25)
  {最优配置 $(N,D,Q)$\\预算变化下的\\配置特征};

\node[archdata] (d) at (-5.8,-7.5)
  {附件 C\\历史模型\\与评测信息};
\node[archmodel] (q4) at (-0.8,-7.5)
  {问题四：能力演进与前沿预测\\损失与评测桥接、回归与外推};
\node[archoutput] (o4) at (4.8,-7.5)
  {规模与非规模贡献\\能力前沿预测区间};

\foreach \archsrc/\archmid/\archdst in {a/q1/o1,b/q2/o2,c/q3/o3,d/q4/o4} {
  \draw[archlink] (\archsrc.east) -- (\archmid.west);
  \draw[archlink] (\archmid.east) -- (\archdst.west);
}
\end{tikzpicture}
\endgroup
